\section{Informe del / de la Director/a}

El estudiante de la Carrera de Ingeniería en Informática: \\
\textbf{Guerrero, Martín}, DNI: 43.404.011, Padrón: 107774 \\
ha cumplido satisfactoriamente con 192 horas de Práctica Profesional Supervisada en 'Empresas de base tecnológica 2' en el marco de su Trabajo Profesional bajo mi supervisión.

El estudiante ha desarrollado satisfactoriamente el Plan de Trabajo oportunamente acordado, por lo que cumplieron con la Práctica Profesional correspondiente, la cual ha servido de base al Trabajo Profesional presentado para su aprobación.

Las tareas asignadas consistieron en participar en el diseño y desarrollo de 'SocioUnido', una plataforma SaaS B2B orientada a modernizar la gestión administrativa, optimizar la recaudación y asegurar el control de accesos en clubes deportivos, asumiendo el rol de responsable de inteligencia artificial, \textit{back-end} y relaciones institucionales. El trabajo del estudiante se destacó por su eficiente implementación de la lógica de negocio del servidor, el innovador desarrollo del asistente virtual con Procesamiento de Lenguaje Natural (PLN) y la sólida configuración de la analítica avanzada y la seguridad lógica del sistema.

El responsable institucional por la organización receptora fue el [Nombre y apellido], en su rol de [Rol en la cátedra] de la materia 'Empresas de base tecnológica 2'. El Mg. Ing. Gabriel Piñeiro, en su carácter de director, estuvo a cargo de la supervisión por parte de la Facultad de Ingeniería.

\vspace{1.5cm}
\noindent\rule{8cm}{0.4pt} \\
FIRMA Y ACLARACIÓN DEL DIRECTOR \\
(completar con firma ológrafa escaneada)

\vspace{1.5cm}
\noindent\rule{8cm}{0.4pt} \\
FIRMA Y ACLARACIÓN DEL RESPONSABLE DE LA ORGANIZACIÓN RECEPTORA \\
(firma ológrafa escaneada)

\vspace{0.5cm}

\section{Informe del Estudiante}

Participar en la creación de 'SocioUnido' desde su concepción hasta convertirse en un Producto Mínimo Viable (MVP) funcional ha sido una experiencia sumamente enriquecedora. Construir un producto desde cero, adoptando el verdadero espíritu emprendedor que promueve la materia 'Empresas de base tecnológica 2', me permitió consolidar mis conocimientos técnicos y aplicarlos a una problemática real del mercado. La oportunidad de desarrollar esta Práctica Profesional Supervisada dentro de este marco académico fue vital, ya que aportó una visión integral sobre el ciclo de vida del \textit{software}, la gestión ágil, el trabajo en equipo y la estrategia de negocio, brindándome herramientas invaluables para mi futuro como profesional de la ingeniería.

\vspace{1.5cm}
\noindent\rule{8cm}{0.4pt} \\
FIRMA Y ACLARACIÓN DEL ESTUDIANTE \\
(completar con firma ológrafa escaneada)